\section{A shorter bound for the opposite-parity Selmer groups}
\label{ms-section}

The entire odd-prime estimate for the opposite-parity summand
admits a classical proof which avoids separating primes dividing
the degree of the field.  Apply the equivariant special-value
theorem over the maximal real subfield: its Galois group then
has order two.  We give the reduction and all numerical estimates.
The classical inputs are Greither--Popescu's cohomological
Coates--Sinnott theorem, Ferrero--Washington's vanishing of the
Iwasawa $\mu$-invariant for abelian number fields, and
Poitou--Tate duality.  These theorems are stated where used;
they are not reproved here.

\subsection{The group and its parity}

Let $E/\mathbb Q$ be an abelian CM extension, $j$ complex
conjugation, $F=E^{\langle j\rangle}$, and $n=[E:\mathbb Q]$.
Choose $N$ with $E\subset\mathbb Q(\zeta_N)$.
Necessarily $N\ge3$, $n$ is even, and $n\le\varphi(N)\le N$.
For an odd prime $p$, let $G_{E,S_p}$ be the Galois group
of the maximal extension unramified away from the places
above $p$ and the infinite places.
The Tate module $\mathbb Z_p(1)$ is the inverse limit of the
$p$-power roots of unity; integer twists are its tensor powers
and their duals.  Set
$\mathbb Q_p/\mathbb Z_p(-e)=
\mathbb Z_p(-e)\otimes_{\mathbb Z_p}\mathbb Q_p/\mathbb Z_p$.
For $e\ge1$ define
\[
 \mathcal S_p(E,e)=
 \ker\left(H^1(G_{E,S_p},\mathbb Q_p/\mathbb Z_p(-e))
       \longrightarrow
       \bigoplus_{v\mid p}H^1(E_v,\mathbb Q_p/\mathbb Z_p(-e))
       \right).
 \tag{MS1}
\]
Cohomology is continuous, and the local maps are restriction.
The action of $j$ conjugates Galois arguments and acts on
the coefficient module.  Since two is a unit in $\mathbb Z_p$,
the two eigenspaces are the images of $(1+j)/2$ and $(1-j)/2$.
We require
\[
 b_p(E,e)=\operatorname{length}_{\mathbb Z_p}
                     \mathcal S_p(E,e)^{\,(-1)^{e+1}}.
 \tag{MS2}
\]
This is the summand called \texttt{ShaMinus} in the source.
It is not always the eigenspace with eigenvalue $-1$.

\begin{theorem}\label{ms-main}
For every integer $t\ge2$ the groups in \textup{(MS2)}, for
$1\le e<t$, are finite and only finitely many are nonzero.
Moreover
\[
 \sum_{\substack{p\ {\rm prime}\\p\ne2}}\log p
       \sum_{e=1}^{t-1} b_p(E,e)
 \le n(t^2+t-2)\log(tN)
 \le \frac98 n t^2\log(tN).
 \tag{MS3}
\]
In particular this supplies the source's estimate
$3nt^2\log(tN)$, including primes dividing $[E:\mathbb Q]$,
without its additional degree-dependent error term.
\end{theorem}

\subsection{The special-value theorem over a group of order two}

For $m\ge2$ put $\epsilon=(-1)^m$ and
$e_\epsilon=(1+\epsilon j)/2$.
Let $S$ consist of the infinite places of $F$ and its
finite places above rational primes dividing $N$.
It contains all places ramified in $E/F$.  Write
\[
 H^i_{S,p}(E,m)=
 H^i_{\mathrm{et}}(\mathcal O_{E,S}[1/p],\mathbb Z_p(m)),
 \qquad
 W_p(E,m)=H^0(E,\mathbb Q_p/\mathbb Z_p(m)).
 \tag{MS4}
\]
The latter group is cyclic of finite order $w_p(E,m)$.
It has $j$-eigenvalue $\epsilon$, because the cyclotomic
character of a complex conjugation is $-1$.

Let $X(E)$ be the group of Dirichlet characters modulo $N$
factoring through $\operatorname{Gal}(E/\mathbb Q)$, and set
\[
 X_\epsilon(E)=\{\chi\in X(E):\chi(-1)=\epsilon\},
 \qquad
 R_m(E,N)=\prod_{\chi\in X_\epsilon(E)} L(1-m,\chi).
 \tag{MS5}
\]
There are $n/2$ factors.  Every character is viewed modulo
$N$, even when its primitive conductor is smaller.
Thus its $L$-function omits the Euler factors at every
prime dividing $N$.

\begin{proposition}[Classical special-value theorem,
specialized]\label{ms-special-value}
For every odd prime $p$,
\[
 \operatorname{length}_{\mathbb Z_p}
                    e_\epsilon H^2_{S,p}(E,m)
       =v_p(w_p(E,m))+v_p(R_m(E,N)).
 \tag{MS6}
\]
The group on the left is finite, and $R_m(E,N)$ is a
nonzero rational number.
\end{proposition}
\begin{proof}
Here is the precise classical input.
For an abelian extension $K/k$ with $k$ totally real,
$p>2$, $m\ge2$, and $S$ containing the infinite and
ramified places, Greither--Popescu
\cite[Theorem 6.11]{greither-popescu} prove
\[
 \begin{split}
 &\operatorname{Ann}_{\mathbb Z_p[G]}
     \bigl(H^1_{\mathrm{et}}(\mathcal O_{K,S}[1/p],
                             \mathbb Z_p(m))_{\mathrm{tors}}\bigr)
                       \Theta_{S,K/k}(1-m)\\
 &\hspace{12mm}=
 e_m(K/k)\operatorname{Fitt}_{\mathbb Z_p[G]}
             H^2_{\mathrm{et}}(\mathcal O_{K,S}[1/p],\mathbb Z_p(m)).
 \end{split}
 \tag{MS7}
\]
The hypothesis is the vanishing of the cyclotomic
$\mu$-invariant for the maximal CM subfield of $K(\mu_p)$.
For $K=E$ that field is abelian over $\mathbb Q$, so
Ferrero--Washington \cite{ct-fw} supplies it.
The idempotent $e_m(K/k)$ is the product, over real places
of $k$, of $(1+(-1)^m j_v)/2$.  For $E/F$ every $j_v$
is $j$, so the product is $e_\epsilon$.
Finiteness of $H^2$ is part of the standard cohomological
setting of this theorem
\cite[\S6.2, preceding Lemma 6.9]{greither-popescu}.

We spell out its specialization.  The maps $j\mapsto1$
and $j\mapsto-1$ give
\[
 \mathbb Z_p[\langle j\rangle]\simeq
                       \mathbb Z_p\times\mathbb Z_p,
 \qquad
 e_\epsilon\mathbb Z_p[\langle j\rangle]\simeq\mathbb Z_p.
 \tag{MS8}
\]
The inverse uses division only by two.  There is no
division by $n$ and no exclusion of $p\mid n$.
The torsion of $H^1_{S,p}(E,m)$ is $W_p(E,m)$
\cite[Lemma 6.9(1)]{greither-popescu}; this also follows
from the exact sequence for
$\mathbb Z_p(m)\subset\mathbb Q_p(m)$, since
$H^0(E,\mathbb Q_p(m))=0$ and $H^1(E,\mathbb Q_p(m))$
is a vector space.
In the $\epsilon$ factor its annihilator is
$(p^{v_p(w_p(E,m))})$.
For a finite $\mathbb Z_p$-module
$M=\bigoplus_i\mathbb Z_p/p^{a_i}$, the diagonal presentation
gives $\operatorname{Fitt}(M)=(p^{\sum_i a_i})$.
The exponent is its length.

The $\epsilon$ component of $\Theta_{S,E/F}(1-m)$ is
$\zeta_{F,S}(1-m)$ for $\epsilon=1$, and
$L_S(1-m,E/F)$ for $\epsilon=-1$.
Artin factorization identifies these with (MS5).
Indeed, induction of the trivial character of
$\langle j\rangle$ to $\operatorname{Gal}(E/\mathbb Q)$
is the sum of the characters with $\chi(j)=1$;
induction of its nontrivial character is the sum with
$\chi(j)=-1$.  Equality of Euler factors proves the
factorization, including the removal of all primes above
the rational divisors of $N$.
Projecting (MS7) to (MS8) proves (MS6).

The character set in (MS5) is invariant under Galois
automorphisms of the field of character values, so the
product is rational.  Every factor is nonzero.
Its primitive character has parity $(-1)^m$; the
functional equation expresses its value at $1-m$ as a
nonzero constant times $L(m,\bar\chi)$.
The absolutely convergent Euler product makes the latter
nonzero.  Each removed Euler factor is
$1-\chi^*(\ell)\ell^{m-1}\ne0$, since $m\ge2$.
\end{proof}

\subsection{Localization and the strict Selmer group}

\begin{lemma}\label{ms-localization}
For $m=e+1\ge2$ and odd $p$,
\[
 b_p(E,e)\le
 \operatorname{length}_{\mathbb Z_p}
                       e_{(-1)^m}H^2_{S,p}(E,m).
 \tag{MS9}
\]
\end{lemma}
\begin{proof}
Apply finite-coefficient Poitou--Tate duality
\cite[Chapter I, Theorem 4.10(a)]{milne-adt}
to $A_a=\mathbb Z/p^a(1-m)$ and its Cartier dual
$T_a=\mathbb Z/p^a(m)$.  It pairs their local restriction
kernels in degrees one and two perfectly.
The pairings respect the inclusions $A_a\longrightarrow A_{a+1}$
and the quotient maps $T_{a+1}\longrightarrow T_a$.
Continuous cohomology with discrete coefficients commutes
with filtered direct limits.  Since there are finitely many
local places, exactness of that limit identifies the direct
limit of the degree-one kernels with (MS1).
All finite-coefficient cohomology groups in question are
finite.  Their inverse systems have the Mittag--Leffler
property (descending sequences of images in a fixed finite
group stabilize); hence their first derived inverse limits
vanish.  Consequently the inverse limit of the degree-two
kernels is the kernel for
$H^2(G_{E,S_p},\mathbb Z_p(m))$.
Pontryagin duality exchanges this inverse limit and the
direct limit just described, giving the required perfect
pairing.  These pairings are equivariant for automorphisms of $E$:
cup products and restriction are natural, and an automorphism
permutes the local invariant maps without changing their sum.
The value group $\mathbb Q_p/\mathbb Z_p$ has trivial
$j$-action.  The eigenspaces of the same sign therefore
pair perfectly; opposite signs pair to zero, since two
is invertible.  Hence $b_p(E,e)$ is the length of a subgroup
of the $(-1)^m$ eigenspace of
$H^2(G_{E,S_p},\mathbb Z_p(m))$.

This is
$H^2_{\mathrm{et}}(\mathcal O_E[1/p],\mathbb Z_p(m))$.
Localization injects it into $H^2_{S,p}(E,m)$.
The potentially preceding term for a removed prime
$v\nmid p$ is
$H^0(\kappa(v),\mathbb Z_p(m-1))$.
Frobenius acts by $(Nv)^{m-1}\ne1$, and multiplication
by $(Nv)^{m-1}-1$ is injective on $\mathbb Z_p$.
That preceding group is zero.
Adding the finitely many primes gives the injection.
It respects $j$, as the set of removed primes is stable.
Projection by $e_{(-1)^m}$ preserves it, proving (MS9)
and finiteness of its left side.
\end{proof}

\subsection{Elementary special-value estimates}

Define the Bernoulli polynomials by
\[
 \frac{z e^{xz}}{e^z-1}
       =\sum_{a\ge0}B_a(x)\frac{z^a}{a!},
 \qquad B_a=B_a(0),
 \tag{MS10}
\]
and put $D_m=m(m+1)!N$ for $m\ge2$.

\begin{lemma}\label{ms-bernoulli}
For every character $\chi$ modulo $N$,
\[
 L(1-m,\chi)=
       -\frac{N^{m-1}}m\sum_{a=1}^N\chi(a)B_m(a/N).
 \tag{MS11}
\]
The number $D_mL(1-m,\chi)$ is an algebraic integer, and
\[
         D_m\le(mN)^m,\qquad |L(1-m,\chi)|\le(mN)^m.
 \tag{MS12}
\]
\end{lemma}
\begin{proof}
For $z>0$, periodicity gives
\[
 \sum_{\nu\ge1}\chi(\nu)e^{-\nu z}
       =\sum_{a=1}^N\chi(a)\frac{e^{-az}}{1-e^{-Nz}}.
\]
Its Mellin integral is $\Gamma(s)L(s,\chi)$ for
$\operatorname{Re}s>1$.
Split that integral at one and subtract sufficiently many
terms of the power series at zero.
The tail at infinity is entire, and a subtracted term
$c z^h$ contributes $c/(s+h)$.
The coefficient of $z^{m-1}$ in the exponential sum is
$(-1)^mN^{m-1}\sum_a\chi(a)B_m(a/N)/m!$, by (MS10).
The residue of $\Gamma(s)$ at $s=1-m$ is
$(-1)^{m-1}/(m-1)!$.
Dividing residues proves (MS11); the same subtraction
supplies the analytic continuation used here.

A denominator bound follows without a theorem about
Bernoulli denominators.  The formal identities
\[
 \frac z{e^z-1}=\sum_{b\ge0}\frac{(-1)^b}{b+1}(e^z-1)^b,
 \qquad
 (e^z-1)^b=b!\sum_{a\ge b}
    \left\{\begin{smallmatrix}a\\b\end{smallmatrix}\right\}
    \frac{z^a}{a!}
\]
give
\[
 B_a=\sum_{b=0}^a\frac{(-1)^b b!}{b+1}
                \left\{\begin{smallmatrix}a\\b\end{smallmatrix}\right\}.
 \tag{MS13}
\]
The braces count partitions into nonempty subsets and
are integers.  For $a\le m$, $(m+1)!$ clears all
denominators in (MS13).
Also $B_m(x)=\sum_{a=0}^m\binom ma B_a x^{m-a}$.
Thus $N^m(m+1)!B_m(b/N)$ is an integer for integer $b$.
Multiplying (MS11) by $D_m$ gives an integral linear
combination of roots of unity and zeros.

Repeated integration by parts in Fourier coefficients,
using $B_a'=aB_{a-1}$ and
$B_a(1)-B_a(0)=0$ for $a\ne1$, gives
\[
 B_m(x)=-\frac{m!}{(2\pi i)^m}
          \sum_{h\ne0}\frac{e^{2\pi ihx}}{h^m},
                  \quad 0\le x\le1,\quad m\ge2.
 \tag{MS14}
\]
The constant coefficient is zero by (MS10).
The series converges absolutely uniformly and has the
same Fourier coefficients as the continuous periodic
function $B_m$, proving equality, including endpoints.
Hence
$|B_m(x)|\le2m!\zeta(m)/(2\pi)^m\le m!$,
using $\zeta(m)\le2$ and $2\pi>4$.
Equation (MS11) gives
$|L(1-m,\chi)|\le N^m(m-1)!\le(mN)^m$.
Finally $m!\le m^{m-1}$ and
$m+1\le3^{m-1}\le N^{m-1}$, so
$m(m+1)!N\le m^m(m+1)N\le(mN)^m$.
\end{proof}

\begin{lemma}\label{ms-roots}
The integer $w_{\mathrm{odd}}(E,m)=\prod_{p\ne2}w_p(E,m)$
is finite and
\[
             \log w_{\mathrm{odd}}(E,m)\le nm\log(mN).
 \tag{MS15}
\]
\end{lemma}
\begin{proof}
If $p^a\mid w_p(E,m)$, the subgroup
$\operatorname{Gal}(E(\zeta_{p^a})/E)$ of
$(\mathbb Z/p^a)^\times$ has exponent dividing $m$.
For odd $p$ it is cyclic, hence has order at most $m$.
Its index is
$[E\cap\mathbb Q(\zeta_{p^a}):\mathbb Q]\le n$.
Therefore $p^{a-1}(p-1)\le nm$ and
$p^a\le3nm/2$.
Only primes $3\le p\le nm+1$ occur.
Since $nm$ is even, there are at most $nm/2$ odd integers
in that interval, whence
\[
 \log w_{\mathrm{odd}}(E,m)
       \le\frac{nm}{2}\log(3nm/2)\le nm\log(mN).
\]
The last inequality uses $n\le N$ and
$3nm/2\le(mN)^2$, since $mN\ge6$.
The prime-power bound also proves finiteness.
\end{proof}

\begin{proof}[Proof of Theorem~\ref{ms-main}]
For fixed $m=e+1$, put $r=n/2$ and $R=R_m(E,N)$.
The number $D_m^rR$ is a nonzero algebraic integer by
Lemma~\ref{ms-bernoulli}, and is rational by
Proposition~\ref{ms-special-value}.  It is therefore a
nonzero rational integer.  So is
$I=w_{\mathrm{odd}}(E,m)D_m^rR$.
Equations (MS6) and (MS9) imply, for odd $p$,
\[
 0\le b_p(E,m-1)\le
       v_p(w_p(E,m))+v_p(R)\le v_p(I).
 \tag{MS16}
\]
The factors $w_q$ for $q\ne p$ are $p$-adic units,
and $D_m^r$ only increases the valuation.
Only primes dividing $I$ can have a nonzero group.
Summing and applying (MS12) and (MS15) yields
\[
 \begin{split}
 \sum_{p\ne2}b_p(E,m-1)\log p
 &\le\log|I|\\
 &\le nm\log(mN)+\frac n2 m\log(mN)+\frac n2 m\log(mN)\\
 &=2nm\log(mN).
 \end{split}
 \tag{MS17}
\]
Using the actual integer $I$ handles possible denominators
of $R$ at two; bounding odd valuations by $\log|R|$ alone
would not be justified.
Sum over $2\le m\le t$ and use
$\log(mN)\le\log(tN)$.
The equality $2\sum_{m=2}^t m=t^2+t-2$ gives the first
inequality in (MS3).
The second is equivalent to $(t-4)^2\ge0$.
\end{proof}

\paragraph{Relation to the submitted proof.}
Taking $N$ to be the conductor proves the required
opposite-parity estimate for every abelian CM field, and
therefore for the composita in the source.
In its notation this supplies \texttt{ThmBcFree} with
the additional constant equal to zero.
This is a separate conventional proof using the named
classical theorems.  It does not certify the source's
individual characterwise annihilator argument.
It does not estimate the complementary eigenspace
$(-1)^e$, whose relevant special value has a zero.


